\chapter{Sampling and Digital Signal Processing}
\label{ch:dsp}
So far, every circuit has been analog: voltage and current exist at every instant.
A modern radio converts its signal to numbers early and processes it in software.
Sampling is still \emph{multiplication}, an operation \cref{sec:mixers} already
analyzed; only the multiplier is new.

This chapter also proves the exact linear phase of a no-feedback digital filter
(\cref{sec:firiir}) and distinguishes the sampling Nyquist theorem from the
Nyquist plot of \cref{sec:nyquist}.

\section{Sampling Is Multiplication}
\label{sec:samplingismult}
To sample a signal is to look at it only at instants spaced \(T_s=1/f_s\) apart. Model
that as multiplying by a train of impulses, one per sampling instant:
\[
x_s(t)=x(t)\cdot p(t),
\qquad
p(t)=\sum_{k=-\infty}^{\infty}\delta(t-kT_s).
\]
Everything follows from asking what \(p(t)\) looks like in the frequency domain. It is
periodic with period \(T_s\), so by \cref{sec:fourier} it is a sum of sinusoids at
multiples of \(f_s\)---and because an impulse train is as sharp in time as a signal can
be, \emph{all} of those harmonics are present, with equal weight:
\[
p(t)=\frac{1}{T_s}\sum_{n=-\infty}^{\infty}\ee^{\jj 2\pi n f_s t}.
\]
That is the key fact: the sampling operation is multiplication by a
\textbf{comb} of frequencies spaced \(f_s\) apart.

Now use the mixer result. \Cref{sec:mixers} showed that multiplying by a sinusoid at
\(f_{\mathrm{LO}}\) shifts a signal to the sum and difference frequencies. Here we are
multiplying by \emph{every} multiple of \(f_s\) at once, so the input spectrum is
shifted to sit around every one of them:
\[
\boxed{\;
X_s(f)=\frac{1}{T_s}\sum_{n=-\infty}^{\infty}X(f-nf_s)\;}
\]
\emph{Sampling replicates the spectrum, once around each multiple of the sampling
rate.} A mixer made two copies because it multiplied by one frequency; a sampler makes
infinitely many because it multiplies by infinitely many. No new mathematics is
required; only the multiplier changes.
This impulse-train model and its replicated spectra are the standard frequency-domain
description of ideal sampling \cite[pp.~23--25]{giannini2008baseband}.

\begin{controlsbox}
\textbf{One phenomenon drives the chapter.} Sampling creates spectrum copies spaced
\(f_s\). The remaining questions are when they overlap (\cref{sec:samplingthm}),
what precision costs (\cref{sec:quantnoise}), and how to change rates without
overlap (\cref{sec:decimation}). Subsequent signal processing is filtering with
poles and zeros placed by arithmetic instead of components.
\end{controlsbox}

\section{The Sampling Theorem, and Aliasing as Folding}
\label{sec:samplingthm}
The copies are harmless as long as they do not overlap. Suppose the signal contains no
frequency above \(f_{\max}\), so each copy of the spectrum occupies a band of width
\(f_{\max}\) either side of its center. The copy at \(0\) reaches up to \(f_{\max}\);
the copy at \(f_s\) reaches down to \(f_s-f_{\max}\). They stay clear of each other
precisely when
\[
f_s-f_{\max}>f_{\max}
\qquad\Longleftrightarrow\qquad
\boxed{\;f_s>2f_{\max}\;}
\]
This is the \textbf{sampling theorem} for a real, strictly band-limited baseband
signal, and it is a non-overlap condition rather than a convention
\cite[pp.~23--25]{giannini2008baseband}. The frequency \(f_s/2\) is the
\textbf{Nyquist frequency}. Equality is an ideal boundary; a realizable
anti-aliasing filter requires transition bandwidth, so practical designs place the
highest wanted component below \(f_s/2\).

What if the condition fails? Then a component above \(f_s/2\) lands inside the
baseband copy, at
\[
f_{\text{alias}}=\bigl|f-nf_s\bigr|
\]
for whichever \(n\) brings it closest to zero. The spectrum \emph{folds} about
\(f_s/2\): a component just above the Nyquist frequency appears just below it, and one
at \(f_s\) appears at DC. That is \textbf{aliasing}.

\begin{mistakebox}
\textbf{Aliasing cannot be undone.} This is the point that matters, and it follows from
the sampling operation rather than from the equipment. Two different input
frequencies---one below \(f_s/2\), one folded down onto it---produce \emph{identical
sample sequences}. Not similar: identical. Once the numbers exist there is no
information left that distinguishes the two, so no filter, algorithm, or amount of
processing applied afterwards can separate them. The only cure is to prevent the
offending component from reaching the converter, which is what an
\textbf{anti-aliasing} low-pass filter is for
\cite[pp.~23--25]{giannini2008baseband}.
\end{mistakebox}

The anti-aliasing filter also explains why a converter's usable bandwidth is
\emph{less} than \(f_s/2\). A real filter does not stop at its cutoff; it rolls off at
\(-20n\,\si{\decibel}\) per decade (\cref{sec:cascadeadd}), so there must be a
transition band between the highest wanted frequency and \(f_s/2\) in which the filter
does its work. Higher filter order narrows the transition band and therefore increases
usable bandwidth. This is the order-versus-sharpness trade of
\cref{sec:filterspecs}, now paid for in sample rate.

\begin{workedbox}[: where does a stray signal land?]
An SDR samples at \(f_s=\SI{48}{\kilo\hertz}\), so its Nyquist frequency is
\SI{24}{\kilo\hertz}. A \SI{30}{\kilo\hertz} spur leaks past the anti-aliasing filter.
With \(n=1\),
\[
f_{\text{alias}}=|30-48|=\SI{18}{\kilo\hertz},
\]
so it appears as an \SI{18}{\kilo\hertz} tone---inside the audio band, sounding
entirely real, and untraceable from the samples alone. Note the folding: the spur was
\SI{6}{\kilo\hertz} \emph{above} Nyquist and its alias sits \SI{6}{\kilo\hertz}
\emph{below} it. A \SI{48}{\kilo\hertz} spur would fold all the way to DC.
\end{workedbox}

\section{Quantization: Six Decibels per Bit}
\label{sec:quantnoise}
Sampling discretizes time. The other half of conversion discretizes \emph{amplitude},
and it has its own cost.

An \(N\)-bit converter spanning a full-scale range \(V_{\mathrm{FS}}\) has
\(2^N\) levels, so its step size is
\[
\Delta=\frac{V_{\mathrm{FS}}}{2^N}.
\]
Rounding to the nearest level makes an error in \(\pm\Delta/2\). For an ideal
converter and a sufficiently varying signal, the usual quantization-noise model
treats that error as uniform and uncorrelated with the input. Its mean-square value
is then the variance of a uniform distribution of width \(\Delta\):
\[
\overline{e^2}=\frac{1}{\Delta}\int_{-\Delta/2}^{\Delta/2}e^2\,\dd e
=\frac{\Delta^2}{12}.
\]
Within those assumptions, quantization behaves as added noise of RMS value
\(\Delta/\sqrt{12}\) \cite[p.~188]{vankka2005synthesizers}. Compare it with a
full-scale sinusoid, whose mean square is \(V_{\mathrm{FS}}^2/8\) by
\cref{sec:rms}:
\[
\mathrm{SNR}=\frac{V_{\mathrm{FS}}^2/8}{\Delta^2/12}
=\frac{12}{8}\,\frac{V_{\mathrm{FS}}^2}{\Delta^2}
=\frac{3}{2}\cdot 2^{2N}.
\]
In decibels, this \emph{ideal full-scale-sine quantization SNR} is
\[
\boxed{\;\mathrm{SNR}_{\dB}=6.02\,N+1.76\;}
\]
This ideal result is derived in \cite[p.~188]{vankka2005synthesizers}.
Thus each additional ideal bit contributes about \SI{6}{\decibel} of
quantization-limited SNR. This is the same power ratio as one nominal S unit
(\cref{sec:sunits}), but it is not by itself a guarantee of a converter's measured
dynamic range.

\begin{workedbox}[: how many bits for a millivolt?]
A signal spans \SI{1}{\volt} and must be resolved to \SI{1}{\milli\volt}. The
requirement is \(2^N\ge1000\), and since \(2^{10}=1024\),
\[
N=10\ \text{bits}
\]
suffices while \(N=9\) (512 levels, \SI{1.95}{\milli\volt} steps) does not. Under
the ideal full-scale-sine model, 10, 16, and 24 bits correspond to approximately
\SI{62}{\decibel}, \SI{98}{\decibel}, and \SI{146}{\decibel} of quantization-limited
SNR. Real converters generally achieve less because thermal noise, nonlinearity,
clock jitter, and other errors reduce their effective number of bits.
\end{workedbox}

\subsection{Oversampling Buys Back Resolution}
Under the ideal white-quantization-noise model, the total noise power
\(\Delta^2/12\) is spread from DC to \(f_s/2\). A receiver listening in a much
narrower bandwidth \(B\), with appropriate digital filtering, admits only the
corresponding fraction, giving
\[
\text{processing gain}=10\log_{10}\frac{f_s/2}{B}\ \si{\decibel}
\]
The corresponding in-band expression appears in
\cite[p.~188]{vankka2005synthesizers}.
Oversampling can therefore improve in-band quantization-limited SNR without adding
nominal bits, although it does not remove thermal noise, aperture jitter, distortion,
or front-end limits. \textbf{Dither} deliberately randomizes signal-correlated
quantization error, trading coherent spurs for a more noise-like spectrum that can
be reduced by filtering or averaging \cite[pp.~113--118]{vankka2005synthesizers}.

\begin{mistakebox}
The two ends of a converter's range fail differently, and this is the practical
difference between an analog and a digital receiver. The noise floor is \emph{soft}: it
can be pushed down by oversampling and filtering. The ceiling is \emph{hard}---once the
input exceeds the reference voltage the converter simply cannot represent it, and the
output clips abruptly rather than compressing gracefully as an overdriven amplifier does
(\cref{sec:compression}). An SDR is certainly overloaded once its converter input
exceeds full scale. In a real receiver, however, front-end compression,
protection-device conduction, or converter nonlinearity can provide an observable
warning region before hard clipping.
\end{mistakebox}

\section{Changing Rate: Decimation and Interpolation}
\label{sec:decimation}
Software radios rarely keep the converter's raw rate. \textbf{Decimation} reduces the
sample rate by keeping every \(M\)th sample; \textbf{interpolation} raises it by
inserting samples. The asymmetry between them is a direct reading of
\cref{sec:samplingthm}.

Decimating divides \(f_s\) by \(M\), which divides the Nyquist frequency by \(M\) too.
Content that was legally below the old \(f_s/2\) may now be above the new one---so it
will fold. The fix is forced: \emph{filter first, then discard samples}. That is why a
decimator always contains an anti-aliasing filter ahead of the rate change, and it is
the same argument as the anti-aliasing filter at the converter input, applied to a rate
that is about to shrink.

Interpolation moves the Nyquist frequency the other way, so nothing that was already
legal can become illegal, and no filter is needed \emph{beforehand}. (One is generally
used afterwards, to remove the images the inserted samples create.)

These anti-aliasing and anti-imaging requirements follow directly from the spectral
replicas created by sampling and rate conversion
\cite[pp.~23--26]{giannini2008baseband}.

\section{FIR and IIR: Proving the Linear-Phase Claim}
\label{sec:firiir}
\Cref{ch:filters} listed two digital filter families and asserted that one of them can
be exactly linear phase. With \cref{sec:groupdelay} in hand that assertion can be
discharged.

A \textbf{finite impulse response} (FIR) filter computes each output as a weighted sum
of the last \(N\) inputs, with no feedback:
\[
y[k]=\sum_{m=0}^{N-1}h_m\,x[k-m].
\]
Its frequency response is \(H(\ee^{\jj\omega T_s})=\sum_m h_m\ee^{-\jj m\omega T_s}\).
Now impose \textbf{symmetry} on the coefficients, \(h_m=h_{N-1-m}\), and pair the terms
from the two ends. Each pair combines into a cosine times a common exponential, and
factoring that exponential out gives
\[
H=\underbrace{\ee^{-\jj\omega T_s(N-1)/2}}_{\text{pure delay}}
\cdot\underbrace{\bigl(\text{a real function of }\omega\bigr)}_{\text{no phase}} .
\]
Away from response zeros, the phase contributed by the first factor is strictly
linear in \(\omega\); the real factor can only add sign changes, seen as phase jumps
at its zeros. Wherever group delay is defined,
\[
\tau_g=-\frac{\dd\phi}{\dd\omega}=\frac{(N-1)T_s}{2}
\]
exactly \((N-1)/2\) samples. Thus a symmetric FIR filter has generalized linear
phase and constant group delay throughout each nonzero response interval. This
proves the claim of \cref{ch:filters}; symmetry, not poles, is sufficient.

An \textbf{infinite impulse response} (IIR) filter feeds its output back:
\[
y[k]=\sum_m b_m x[k-m]-\sum_{j\ge1}a_j\,y[k-j].
\]
Feedback introduces poles. General IIR designs therefore have phase shaped by both
poles and zeros and do not automatically possess the exact generalized-linear-phase
property above. For many conventional magnitude specifications, an IIR realization
uses fewer arithmetic operations than a comparably selective FIR, but that is a
design trade rather than a universal theorem.

\begin{controlsbox}
\textbf{FIR versus IIR.} An FIR filter is a finite delay cascade with no recursive
loop. Its response to a finite-duration input eventually ends, and coefficient
symmetry can give exact generalized linear phase, but high selectivity may require
many taps. An IIR filter uses feedback and can produce a decaying response of
unbounded duration; it is unstable if a pole lies outside the unit circle.
\end{controlsbox}

\section{Getting Back Out: the Zero-Order Hold}
A converter running the other way holds each sample's value until the next one arrives,
producing a staircase rather than a smooth curve. That hold is itself a filter: holding
for \(T_s\) is convolution with a rectangular pulse of width \(T_s\), whose frequency
response is
\[
H_{\mathrm{ZOH}}(f)=T_s\,\frac{\sin\pi fT_s}{\pi fT_s},
\]
a \(\operatorname{sinc}\) shape \cite[pp.~186--188]{vankka2005synthesizers}.
It \emph{droops}: the response is already down about \SI{3.9}{\decibel} at
\(f_s/2\), a known and correctable error. Its nulls lie at multiples of \(f_s\),
which only partially suppresses the staircase's spectral images. A
\textbf{reconstruction}, or anti-imaging, low-pass filter therefore follows the
converter. The input anti-aliasing filter prevents out-of-band content from folding
into baseband; the output reconstruction filter removes spectral images created by
conversion.

\section{Direct Digital Synthesis}
The same ideas run backwards make an oscillator. A \textbf{direct digital synthesizer}
adds a constant \(M\) to an \(N\)-bit accumulator on every clock edge; the accumulator
overflows periodically, and that periodicity is the output frequency:
\[
f_{\mathrm{out}}=f_{\mathrm{clk}}\,\frac{M}{2^N}.
\]
The accumulator architecture and tuning relation are developed in
\cite[pp.~61--63]{vankka2005synthesizers}.
Because \(f_{\mathrm{clk}}\) can come from a crystal, the output inherits crystal
stability while remaining tunable in steps of \(f_{\mathrm{clk}}/2^N\)---which is why
DDS replaced the variable-frequency oscillator. It is the digital alternative to the
phase-locked synthesizer of \cref{sec:pll}: a DDS sets frequency by arithmetic and a
PLL by a feedback loop, and the two are often combined, a PLL multiplying a
DDS-generated reference.

The accumulator produces phase, so a phase-to-amplitude table supplies waveform
samples to a converter. Finite tuning-word width limits the set of exactly
representable frequencies. Discrete spurs arise more directly from periodic phase
truncation, finite table/amplitude resolution, and DAC nonlinearity
\cite[pp.~97--105]{vankka2005synthesizers}.

\section{Worked Example}
A direct-sampling receiver digitizes the whole HF spectrum with a 16-bit converter
running at \(f_s=\SI{122.88}{\mega\hertz}\), then filters down to a \SI{2.4}{\kilo\hertz}
SSB channel.

\emph{Usable bandwidth.} The Nyquist frequency is \SI{61.44}{\mega\hertz}, so with a
reasonable anti-aliasing filter the receiver covers HF to \SI{30}{\mega\hertz}
comfortably---the transition band from \SI{30}{\mega\hertz} to \SI{61.44}{\mega\hertz}
is more than an octave, which even a modest filter order can fill.

\emph{Ideal converter SQNR.} From \cref{sec:quantnoise},
\[
\mathrm{SNR}_{\rm ideal}=6.02(16)+1.76\approx\SI{98}{\decibel}
\]
across the full \SI{61.44}{\mega\hertz} Nyquist band.

\emph{Ideal in-band quantization-limited SQNR.} If quantization noise is white and
other noise and distortion sources are neglected, the processing gain is
\[
10\log_{10}\frac{61.44\times10^{6}}{2.4\times10^{3}}
=10\log_{10}(2.56\times10^{4})\approx\SI{44}{\decibel},
\]
so this ideal model predicts about \(98+44=\SI{142}{\decibel}\) in the selected
channel \cite[p.~188]{vankka2005synthesizers}. The result demonstrates why a high
sample-rate-to-channel-bandwidth ratio is valuable, not what the complete receiver
must achieve. Effective number of bits, clock jitter and phase noise, thermal noise,
front-end linearity, and strong out-of-band signals set the realizable system range.

\begin{exambox}
Five facts carry most of the questions here, and each is a line of this chapter.
Sample at more than \emph{twice} the highest frequency, or components fold about
\(f_s/2\) and are unrecoverable. An anti-aliasing filter must come \emph{before} the
converter, and before any decimation, but is not needed before interpolation.
Resolution is \(2^N\) levels, so an 8-bit converter has 256 and each bit is worth
\SI{6}{\decibel}. A converter clips once its input exceeds full scale, although a
practical receiver's analog front end may compress first. And an FIR filter can be
exactly linear phase because symmetric coefficients make its group delay constant.
\end{exambox}
